\documentclass[11pt]{article}
\usepackage[margin=1in]{geometry}
\usepackage{amsmath, amssymb}

\title{Exercise 3.13: $q_\pi$ in Terms of $v_\pi$ and $p$}
\author{Sutton \& Barto, Section 3.5}
\date{}

\begin{document}
\maketitle

\section*{Problem}

Give an equation for $q_\pi$ in terms of $v_\pi$ and the four-argument
dynamics function $p$.

\section*{Solution}

The action-value $q_\pi(s,a)$ is the expected return after taking action $a$
in state $s$ and then following $\pi$. After taking the action, the environment
transitions to some next state $s'$ with reward $r$ according to $p(s',r \mid s,a)$.
From $s'$ onward, the expected return is $v_\pi(s')$.

Using $G_t = R_{t+1} + \gamma\, G_{t+1}$:

\begin{align*}
    q_\pi(s, a) &= \mathbb{E}_\pi[R_{t+1} + \gamma\, G_{t+1} \mid S_t = s, A_t = a] \\
                &= \sum_{s'} \sum_{r} p(s', r \mid s, a) \left[ r + \gamma\, v_\pi(s') \right]
\end{align*}

Therefore:
\[
    \boxed{q_\pi(s, a) = \sum_{s'} \sum_{r} p(s', r \mid s, a) \left[ r + \gamma\, v_\pi(s') \right]}
\]

where $s' \in \mathcal{S}^+$ and $r \in \mathcal{R}$.

\subsection*{Connection to the Bellman equation}

Substituting the result from Exercise~3.12 ($v_\pi(s) = \sum_a \pi(a|s)\, q_\pi(s,a)$)
into this equation, or vice versa, yields the Bellman equation~(3.14):
\[
    v_\pi(s) = \sum_{a} \pi(a \mid s) \sum_{s'} \sum_{r}
    p(s', r \mid s, a) \left[ r + \gamma\, v_\pi(s') \right]
\]

Exercises 3.12 and 3.13 decompose the Bellman equation into its two constituent
parts, corresponding to the two levels of the backup diagram (Figure~3.4).

\end{document}
